\chapter{Per-Unit-System und Impedanzumrechnung}
\label{chap:perunit}

\begin{lernziel}
Nach diesem Kapitel können Sie Bezugsgrößen eines Drehstromnetzes festlegen,
Impedanzen auf eine andere Spannungsebene beziehen, aus der
Kurzschlussspannung eines Transformators dessen Impedanz bestimmen und
mehrere Netzkomponenten auf eine gemeinsame Per-Unit-Basis umrechnen. Sie
können außerdem die Grenzen einer überschlägigen Rechnung benennen.
\end{lernziel}

\section{Warum Bezugsgrößen helfen}

Ein Netzmodell kann eine Einspeisung, mehrere Transformatoren und Leitungen
verschiedener Spannungsebenen enthalten. In einem rein ohmsch-induktiven
Ersatzschaltbild müssen alle Impedanzen, die addiert werden, auf dieselbe
Spannungsebene bezogen sein. Das klassische Umrechnen über jedes
Transformatorübersetzungsverhältnis ist korrekt, wird bei mehreren Ebenen aber
unübersichtlich. Das Per-Unit-System (p.u.) beschreibt eine Größe als Anteil
einer festgelegten Bezugsgröße.

Für ein symmetrisches Drehstromsystem legen wir üblicherweise eine
Bezugs-Scheinleistung $S_B$ und eine Leiter-Leiter-Bezugsspannung $U_B$ fest.
Daraus ergeben sich die Bezugsgrößen
\begin{align}
  I_B &= \frac{S_B}{\sqrt{3}\,U_B},\\
  Z_B &= \frac{U_B^2}{S_B},\\
  Y_B &= \frac{1}{Z_B}=\frac{S_B}{U_B^2}.
\end{align}
Dabei sind $S_B$ die gesamte dreiphasige Scheinleistung und $U_B$ die
Leiter-Leiter-Spannung. Mit $U_B$ in Volt und $S_B$ in Voltampere erhält man
$Z_B$ in Ohm. In der Praxis dürfen konsistente Präfixe wie kV und MVA
verwendet werden, sofern das Ergebnis dimensionsrichtig gebildet wird.

Die normierte Impedanz lautet
\begin{equation}
  z_{\mathrm{pu}}=\frac{Z}{Z_B}=Z\frac{S_B}{U_B^2}.
  \label{eq:zpu}
\end{equation}
Analog gilt $i_{\mathrm{pu}}=I/I_B$ und $u_{\mathrm{pu}}=U/U_B$.
Bei Wechselstromrechnung sind Impedanzen und Spannungen komplexe Größen; die
Bezugsgrößen selbst sind reelle positive Beträge.

\begin{merke}
Die Formel $Z_B=U_B^2/S_B$ gilt in dieser Form für ein konsistentes
Drehstromsystem mit Leiter-Leiter-Bezugsspannung und dreiphasiger
Bezugsleistung. Vermischen Sie diese Größen nicht mit einer einphasigen
Rechnung, ohne die Definitionen entsprechend anzupassen.
\end{merke}

\section{Impedanz auf eine andere Spannungsebene beziehen}

Für einen idealen Transformator mit Übersetzungsverhältnis
\begin{equation}
  a=\frac{U_1}{U_2}
\end{equation}
folgt aus der Erhaltung der Leistung und dem Übersetzungsverhältnis der
Spannungen die auf Seite 1 bezogene Impedanz einer Komponente auf Seite 2:
\begin{equation}
  Z_2'=a^2 Z_2.
  \label{eq:impedanz_beziehen}
\end{equation}
In Gegenrichtung gilt $Z_1'=Z_1/a^2$. Die Regel gilt für die gesamte komplexe
Impedanz, also sowohl für den Widerstands- als auch für den Reaktanzanteil:
\begin{equation}
  R'=a^2R,\qquad X'=a^2X.
\end{equation}

\begin{rechenbeispiel}
Eine Niederspannungs-Komponente habe $Z_2=(0{,}08+j0{,}06)\,\Omega$. Ein
idealer Transformator setze $20\,\mathrm{kV}$ auf $0{,}4\,\mathrm{kV}$ um.
Dann ist $a=50$ und damit
\[
 Z_2'=50^2(0{,}08+j0{,}06)\,\Omega
     =(200+j150)\,\Omega.
\]
Der Betrag der Impedanz wird mit $a^2=2500$ multipliziert. Das ist kein
physikalisches Anwachsen der Komponente, sondern die Darstellung derselben
Wirkung auf der Hochspannungsseite des Ersatzschaltbilds.
\end{rechenbeispiel}

\section{Per-Unit-Umrechnung zwischen Bezugsbasen}

Eine physikalische Impedanz bleibt gleich, ihr p.u.-Wert hängt aber von der
gewählten Bezugsleistung und Bezugsspannung ab. Ist eine Impedanz auf der
alten Basis $(S_{B,\mathrm{alt}},U_{B,\mathrm{alt}})$ als
$z_{\mathrm{pu,alt}}$ bekannt, ergibt sich auf einer neuen Basis
\begin{equation}
 z_{\mathrm{pu,neu}}=z_{\mathrm{pu,alt}}
 \frac{S_{B,\mathrm{neu}}}{S_{B,\mathrm{alt}}}
 \left(\frac{U_{B,\mathrm{alt}}}{U_{B,\mathrm{neu}}}\right)^2.
 \label{eq:pu_basiswechsel}
\end{equation}
Die Formel folgt direkt daraus, dass $Z=z_{\mathrm{pu}}U_B^2/S_B$ auf beiden
Basen dieselbe reale Impedanz beschreiben muss.

Wählt man die Bezugsspannungen auf beiden Seiten eines Transformators im
Verhältnis seiner Nennspannungen, vereinfacht sich die Umrechnung über den
Transformator hinweg: Die p.u.-Impedanz bleibt bei gleicher Bezugsleistung
auf beiden Seiten gleich. Genau das macht die Methode für mehrstufige Netze
so nützlich.

\begin{merke}
Beim Basiswechsel darf der Spannungsfaktor in
\cref{eq:pu_basiswechsel} nicht weggelassen werden, wenn die alte und neue
Bezugsspannung nicht identisch sind. Die häufig verwendete Kurzform
$z_{\mathrm{pu,neu}}=z_{\mathrm{pu,alt}}S_{B,\mathrm{neu}}/S_{B,\mathrm{alt}}$
setzt voraus, dass die Bezugsspannung unverändert bleibt oder dass die
Bezugsbasen entlang eines Transformators passend gewählt sind.
\end{merke}

\section{Transformatorimpedanz aus Typenschilddaten}

Auf dem Transformator-Typenschild findet sich häufig die relative
Kurzschlussspannung $u_k$ in Prozent. Sie ist die Spannung, die bei
kurzgeschlossener Gegenseite erforderlich ist, damit bei Nennfrequenz der
Nennstrom fließt, angegeben als Anteil der Nennspannung. Für das übliche
vereinfachte Ersatzschaltbild ergibt sich der Betrag der Impedanz bezogen auf
die Nennbasis des Transformators zu
\begin{equation}
  z_{T,\mathrm{pu}}\approx\frac{u_k[\%]}{100}.
  \label{eq:uk_pu}
\end{equation}
Die Näherung bezieht sich auf den Impedanzbetrag. $u_k$ allein liefert nicht
die Aufteilung in Widerstand und Reaktanz. Wenn die Kupferverluste $P_k$
(Nennstrom-Kurzschlussverluste) bekannt sind, lässt sich diese Aufteilung
näherungsweise ergänzen.

Für die Bezugsleistung $S_N$ und die Leiter-Leiter-Nennspannung $U_N$ ist der
Nennstrom
\[
 I_N=\frac{S_N}{\sqrt{3}U_N}.
\]
Der Impedanzbetrag auf der jeweiligen Seite beträgt
\begin{equation}
 |Z_T|=\frac{u_k}{100}\frac{U_N^2}{S_N},
 \label{eq:zt_ohm}
\end{equation}
sofern $u_k$ als dimensionsloser Wert eingesetzt wird. Für die Umrechnung auf
eine andere Bezugsbasis wird anschließend
\cref{eq:pu_basiswechsel} verwendet.

\begin{rechenbeispiel}
Ein Transformator hat $S_N=\SI{10}{\mega\volt\ampere}$,
$U_N=\SI{20}{\kilo\volt}$ auf der betrachteten Seite und $u_k=\SI{8}{\percent}$.
Dann gilt
\[
 Z_B=\frac{(20\,\mathrm{kV})^2}{10\,\mathrm{MVA}}=40\,\Omega,
 \qquad |Z_T|\approx0{,}08\cdot40\,\Omega=3{,}2\,\Omega.
\]
Das ist der Betrag der auf diese Seite bezogenen Kurzschlussimpedanz, nicht
automatisch ein reiner Blindwiderstand von $j3{,}2\,\Omega$.
\end{rechenbeispiel}

\section{Netzimpedanz aus Kurzschlussleistung}

Ist am Einspeisepunkt die dreiphasige Anfangs-Kurzschlussleistung
$S_k''$ bekannt, kann der Betrag der äquivalenten Netzimpedanz überschlägig
bestimmt werden:
\begin{equation}
 |Z_Q|\approx\frac{U_n^2}{S_k''}.
 \label{eq:zquelle}
\end{equation}
Auf einer p.u.-Basis mit $U_B=U_n$ folgt unmittelbar
\begin{equation}
 |z_{Q,\mathrm{pu}}|\approx\frac{S_B}{S_k''}.
 \label{eq:zquellepu}
\end{equation}
Die Kurzschlussleistung allein bestimmt wiederum nicht den Winkel der
Impedanz. Für eine belastbare Kurzschlussrechnung werden, soweit erforderlich,
auch das Verhältnis $R/X$, die Netzbetriebsbedingungen, die Beiträge von
Motoren oder dezentralen Einspeisern und die anzuwendenden Rechenregeln
berücksichtigt.

\section{Durchgängiges Beispiel: Quelle, Transformator und Kabel}

Gegeben sei ein vereinfachtes Netz mit folgenden Daten:
\begin{itemize}
 \item Einspeisung auf der Mittelspannungsseite: $U_n=\SI{20}{\kilo\volt}$,
       $S_k''=\SI{250}{\mega\volt\ampere}$,
 \item Transformator: $\SI{20}{\kilo\volt}/\SI{0.4}{\kilo\volt}$,
       $S_N=\SI{1}{\mega\volt\ampere}$, $u_k=\SI{6}{\percent}$,
 \item Niederspannungskabel: $Z_K=(0{,}015+j0{,}010)\,\Omega$,
       auf die betrachtete Strecke bezogen.
\end{itemize}

Wir wählen $S_B=\SI{1}{\mega\volt\ampere}$ als gemeinsame
Bezugsleistung. Die Bezugsspannungen sind $\SI{20}{\kilo\volt}$ auf der
Mittelspannungsseite und $\SI{0.4}{\kilo\volt}$ auf der
Niederspannungsseite. Sie stehen im selben Verhältnis wie die
Transformator-Nennspannungen.

\paragraph{1. Einspeisung.}
Mit \cref{eq:zquellepu} folgt
\[
 |z_{Q,\mathrm{pu}}|\approx\frac{1}{250}=0{,}004.
\]
Für die Beispielrechnung nehmen wir vereinfachend eine rein induktive
Quellenimpedanz $z_Q\approx j0{,}004$ an. Das ist eine Modellannahme, keine
Information, die sich allein aus $S_k''$ ableiten lässt.

\paragraph{2. Transformator.}
Da die Bezugsleistung gleich der Transformator-Nennleistung ist und die
Bezugsspannungen den Nennspannungen folgen, gilt
\[
 z_{T,\mathrm{pu}}\approx j0{,}06,
\]
unter der zusätzlichen vereinfachenden Annahme, dass die Transformatorimpedanz
vollständig reaktiv ist. Aus $u_k$ allein wäre diese Annahme nicht beweisbar.

\paragraph{3. Kabel.}
Auf der Niederspannungsseite ist
\[
 Z_{B,\mathrm{NS}}=\frac{(0{,}4\,\mathrm{kV})^2}{1\,\mathrm{MVA}}
 =0{,}16\,\Omega.
\]
Somit
\[
 z_{K,\mathrm{pu}}=\frac{0{,}015+j0{,}010}{0{,}16}
 =0{,}09375+j0{,}0625.
\]

\paragraph{4. Gesamtimpedanz und Fehlerstrom.}
Alle Impedanzen liegen nun auf derselben p.u.-Basis und können addiert werden:
\begin{align*}
 z_{\Sigma,\mathrm{pu}}&=z_Q+z_T+z_K\\
 &\approx0{,}09375+j(0{,}004+0{,}06+0{,}0625)\\
 &=0{,}09375+j0{,}1265.
\end{align*}
Der Betrag ist $|z_{\Sigma,\mathrm{pu}}|\approx0{,}1574$. Der Bezugsstrom auf
der Niederspannungsseite beträgt
\[
 I_{B,\mathrm{NS}}=\frac{1\,\mathrm{MVA}}{\sqrt{3}\cdot0{,}4\,\mathrm{kV}}
 \approx\SI{1.443}{\kilo\ampere}.
\]
Für einen dreipoligen Fehler am Kabelende ergibt die vereinfachte Rechnung
\[
 I_k\approx\frac{I_B}{|z_{\Sigma,\mathrm{pu}}|}
 \approx\SI{9.17}{\kilo\ampere}.
\]

\begin{sicherheitshinweis}
Dieses Beispiel illustriert die Rechenmethode, ist aber keine vollständige
Kurzschlussberechnung nach einer anzuwendenden Norm. Es vernachlässigt unter
anderem den ohmschen Anteil der Quellen- und Transformatorimpedanz,
Spannungsfaktoren, Toleranzen, Betriebszustände und mögliche weitere
Kurzschlussstrombeiträge. Für die Auslegung von Schaltgeräten, Schutzorganen
und Betriebsmitteln müssen die maßgeblichen Randbedingungen und aktuellen
Regelwerke berücksichtigt werden.
\end{sicherheitshinweis}

\section{Typische Fehler bei der Umrechnung}
\begin{itemize}
 \item Leiter-Leiter-Spannung und Phasenspannung werden in einer Rechnung
       vermischt.
 \item Eine Impedanz wird ohne Quadrat des Übersetzungsverhältnisses auf eine
       andere Spannungsebene bezogen.
 \item Bei einem Basiswechsel wird der Faktor für die Bezugsleistung oder
       Bezugsspannung vergessen.
 \item Der Prozentwert $u_k=6\,\%$ wird als $6$ statt als $0{,}06$ eingesetzt.
 \item Aus $u_k$ oder $S_k''$ wird fälschlich ein vollständiger komplexer
       Impedanzwert einschließlich Winkel abgeleitet.
 \item Impedanzen in Ohm von unterschiedlichen Spannungsebenen werden direkt
       addiert, ohne sie vorher umzurechnen.
\end{itemize}

\begin{aufgaben}
 \item Bestimmen Sie $Z_B$ für $S_B=\SI{5}{\mega\volt\ampere}$ und
       $U_B=\SI{10}{\kilo\volt}$.
 \item Ein Transformator hat $S_N=\SI{5}{\mega\volt\ampere}$,
       $U_N=\SI{10}{\kilo\volt}$ und $u_k=\SI{7}{\percent}$. Bestimmen Sie
       den Impedanzbetrag auf dieser Seite in Ohm.
 \item Eine Impedanz beträgt auf einer Basis $z_{\mathrm{pu,alt}}=0{,}08$.
       Die Bezugsspannung bleibt gleich; die Bezugsleistung steigt von
       $\SI{10}{\mega\volt\ampere}$ auf $\SI{20}{\mega\volt\ampere}$.
       Wie groß ist $z_{\mathrm{pu,neu}}$?
 \item Warum reicht die Angabe $u_k=\SI{6}{\percent}$ nicht aus, um den
       vollständigen komplexen Transformatorimpedanzwert zu bestimmen?
 \item Welche Bedingung muss erfüllt sein, bevor Impedanzen in Ohm direkt
       addiert werden dürfen?
\end{aufgaben}

\section*{Kurzlösungen}
\begin{enumerate}[label=\arabic*.]
 \item $Z_B=(10\,\mathrm{kV})^2/(5\,\mathrm{MVA})=20\,\Omega$.
 \item $Z_B=20\,\Omega$ und $|Z_T|\approx0{,}07\cdot20=1{,}4\,\Omega$.
 \item $z_{\mathrm{pu,neu}}=0{,}08\cdot20/10=0{,}16$.
 \item $u_k$ liefert den Impedanzbetrag; für den Winkel bzw. die Aufteilung in
       $R$ und $X$ braucht man weitere Angaben, etwa Kurzschlussverluste.
 \item Sie müssen auf dieselbe Spannungsebene bezogen sein und sich auf
       kompatible Ersatzschaltbild- und Frequenzannahmen beziehen.
\end{enumerate}
